% from main.tex
%%%%%%%%%%  scan idx 565  |  folio 552  %%%%%%%%%%
\begin{equation}\tag{6.1.1}
 M^n,\qquad \bigotimes^n M,\qquad \bigwedge^n M,\quad\text{etc.}
\end{equation}
Les
foncteurs \((\ )^n\), \(\bigotimes\limits^n\) et \(\bigwedge\limits^n\) transforment modules libres (resp. localement libres) en modules libres (resp. localement libres).

De façon générale, considérons des sections \(n_1,\ldots,n_r\) de \(\mathbb Z_X\) et un multifoncteur (covariant, pour fixer les idées)
\[
 T:\Mod^{n_1}(\mathcal O)\times\cdots\times
 \Mod^{n_r}(\mathcal O)\longrightarrow\Mod(\mathcal O)
\]
à \(r\) variables, où \(\Mod_{\mathcal O}^{n}(X)\) désigne la catégorie fibrée sur \(X\) des \(\mathcal O\)-modules localement libres de rang \(n\). Soit
\[
 M=T(\mathcal O^{n_1},\ldots,\mathcal O^{n_r})
\]
et
\[
 G=\operatorname{G\ell}(n_1)\times\cdots\times\operatorname{G\ell}(n_r),
 \quad\text{où }\operatorname{G\ell}(n_i)=\Aut_{\mathcal O}(\mathcal O^{n_i}).
\]
Il s'ensuit immédiatement que \(T\) définit une représentation
\begin{equation}\tag{6.1.2}
 \alpha_T:G\longrightarrow\Aut_{\mathcal O}(M).
\end{equation}
Il est clair que si \(M\) est lui-même localement libre de rang \(m\), pour tout foncteur composé,
\[
 T'T,\ \text{où } T':\Mod^{m}(\mathcal O)\longrightarrow\Mod(\mathcal O),
\]
on a
\begin{equation}\tag{6.1.3}
 \alpha_{T'T}=\alpha_{T'}\,\alpha_T.
\end{equation}

Réciproquement, étant donné la représentation \(\alpha=\alpha_T:G\to\Aut_{\mathcal O}(M)\), on récupère \(T(E_1,\ldots,E_r)\) (pour toute suite \(E_1,\ldots,E_r\) de \(\mathcal O\)-
